\documentclass[a4paper, 14pt]{extarticle}

\usepackage{cmap}
\usepackage[T2A]{fontenc}
\usepackage[utf8]{inputenc}
\usepackage[english, russian]{babel}
\usepackage{amsmath}
\usepackage{subcaption}
\usepackage{graphicx}
\usepackage{multirow}
\usepackage{float}
\usepackage{placeins}
\usepackage{hyperref}

\usepackage[left=3cm, right=1.5cm, top=2cm, bottom=2cm]{geometry}
\usepackage{setspace}
\onehalfspacing 

\usepackage{indentfirst}
\setlength{\parindent}{1.25cm} 
\setlength{\emergencystretch}{3em}
\usepackage[expansion=false,protrusion=false,nopatch=footnote]{microtype}

\usepackage{tocloft}
\addto\captionsrussian{\renewcommand{\contentsname}{СОДЕРЖАНИЕ}}
\renewcommand{\cfttoctitlefont}{\hfill\normalfont}
\renewcommand{\cftaftertoctitle}{\hfill\vspace{\baselineskip}}
\renewcommand{\cftsecdotsep}{\cftdotsep}

\begin{document}

\begin{titlepage}
    \begin{center}
        Министерство науки и высшего образования Российской Федерации\\
        Федеральное государственное автономное образовательное учреждение высшего образования\\
        «Санкт-Петербургский политехнический университет Петра Великого»\\
        
        \textbf{Физико-механический институт}\\
        \textbf{Высшая школа прикладной математики и вычислительной физики}
    \end{center}

    \vfill
    
    \begin{center}
        \textbf{Отчёт по лабораторной работе \textnumero~ 5\\[0.5cm]}
        \textbf{Корреляционный анализ и эллипсы рассеивания\\}
        по дисциплине <<Математическая статистика>>
    \end{center}

    \vfill

    \hfill
    \begin{minipage}{0.45\textwidth}
        \textbf{Выполнил студент гр. 5030102/30202:}\\
        Шильдт Д.С.\\[0.5cm]
        \textbf{Преподаватель:}\\ 
        Баженов А.Н.
    \end{minipage}

    \vfill

    \begin{center}
        Санкт-Петербург\\2026
    \end{center}
\end{titlepage}

\setcounter{page}{2}

\tableofcontents
\newpage

\section{Постановка задачи}

\textbf{Задание:}
\begin{enumerate}
    \item Сгенерировать двумерные выборки $(x, y)$ объёмом $n = 20, 60, 100$ из:
    \begin{itemize}
        \item Двумерного нормального распределения $N(0,0,1,1,\rho)$ с $\rho = 0,0.5,0.9$.
        \item Смеси: $0.9N(0,0,1,1,0.9) + 0.1N(0,0,10,10,-0.9)$.
    \end{itemize}
    \item Для каждой выборки вычислить коэффициенты корреляции Пирсона, Спирмена и квадрантный. Повторить 1000 раз, найти средние и дисперсии.
    \item Построить диаграммы рассеяния и эллипсы равновероятности.
\end{enumerate}

\textbf{Цели и задачи:}
\begin{itemize}
    \item Сравнить исходное значение коэффициента корреляции с вычисленными.
    \item Изучить способы визуализации ковариационной структуры данных.
\end{itemize}

\section{Теоретическая часть}

Двумерная случайная величина $(X,Y)$ называется распределённой нормально, если её плотность вероятности определена формулой:
\begin{equation*}
\begin{split}
&N(x, y, \overline{x}, \overline{y}, \sigma_x, \sigma_y, \rho) = \frac{1}{2\pi\sigma_x\sigma_y\sqrt{1-\rho^2}} \times \\
&\qquad \times \exp \left\{ -\frac{1}{2(1-\rho^2)} \left[ \frac{(x-\overline{x})^2}{\sigma_x^2} - 2\rho\frac{(x-\overline{x})(y-\overline{y})}{\sigma_x\sigma_y} + \frac{(y-\overline{y})^2}{\sigma_y^2} \right] \right\}
\end{split}
\end{equation*}
где $\overline{x},\overline{y}$ --- математические ожидания, $\sigma_x, \sigma_y$ — средние квадратические отклонения компонент, $\rho$ — коэффициент корреляции.

Для количественной оценки степени линейной зависимости между компонентами выборки $\{(x_i, y_i)\}_{i=1}^n$ применяется выборочный коэффициент корреляции Пирсона $r$:
\begin{equation*}
    r = \frac{\frac{1}{n}\sum_{i=1}^n (x_i - \bar{x})(y_i - \bar{y})}{\sqrt{\frac{1}{n}\sum_{i=1}^n (x_i - \bar{x})^2 \sum_{i=1}^n (y_i - \bar{y})^2}} = \frac{K}{s_X s_Y},
\end{equation*}
где $K,s^2_X,s^2_Y$ --- выборочные ковариации и дисперсии с.в. X и Y.

Выборочный коэффициент ранговой корреляции Спирмена определяется как выборочный коэффициент корреляции Пирсона между рангами $u, v$ переменных $X,Y$:
\begin{equation*}
    r_S = \frac{\frac{1}{n} \sum(u_i - \overline{u})(v_i - \overline{v})}{\sqrt{\frac{1}{n} \sum(u_i - \overline{u})^2 \frac{1}{n} \sum(v_i - \overline{v})^2}},
\end{equation*}
где $\overline{u} = \overline{v} = \frac{1 + 2 + \dots + n}{n} = \frac{n + 1}{2}$ --- среднее значение рангов.

Выборочный квадрантный коэффициент корреляции:
\begin{equation*}
    r_Q = \frac{(n_1 + n_3) - (n_2 + n_4)}{n}.
\end{equation*}
где $n_1, n_2, n_3$ и $n_4$ --- количества точек с координатами $(x_i,y_i)$, попавших соответственно в 1, 2, 3 и 4 квадранты декартовой системы с осями $x' = x - \mathrm{med}(x)$, $y' = y - \mathrm{med}(y)$ и с центром в точке $(\mathrm{med}(x), \mathrm{med}(y))$.

Для смеси распределений в данной работе используется линейная комбинация двух независимых выборок:
\begin{equation*}
    X = 0.9X_1 + 0.1X_2, \qquad Y = 0.9Y_1 + 0.1Y_2,
\end{equation*}
где $(X_1, Y_1) \sim N(0,0,1,1,0.9)$, $(X_2, Y_2) \sim N(0,0,10,10,-0.9)$. Так как компоненты независимы и имеют нулевые математические ожидания, для смеси получаем
\begin{equation*}
    D_X = D_Y = 0.9^2 \cdot 1 + 0.1^2 \cdot 10 = 0.91,
\end{equation*}
\begin{equation*}
    \mathrm{cov}(X,Y) = 0.9^2 \cdot 0.9 + 0.1^2 \cdot (-9) = 0.639.
\end{equation*}
Следовательно, коэффициент корреляции смеси равен
\begin{equation*}
    r_{XY} = \frac{\mathrm{cov}(X,Y)}{\sqrt{D_X D_Y}} = \frac{0.639}{\sqrt{0.91 \cdot 0.91}} \approx 0.7022.
\end{equation*}

Геометрическим местом точек плоскости, в которых плотность двумерного нормального распределения сохраняет постоянное значение, выступают эллипсы равновероятности (эллипсы рассеивания). Аналитическое уравнение проекции эллипса:
\begin{equation*}
    \frac{(x - \bar{x})^2}{\sigma_x^2} - 2\rho\frac{(x - \bar{x})(y - \bar{y})}{\sigma_x\sigma_y} + \frac{(y - \bar{y})^2}{\sigma_y^2} = const.
\end{equation*}
Ориентация эллипса относительно координатных осей жестко зависит от коэффициента корреляции $\rho$. Угол наклона главных осей симметрии $\alpha$ вычисляется из соотношения:
\begin{equation*}
    \tan 2\alpha = \frac{2\rho\sigma_x\sigma_y}{\sigma_x^2 - \sigma_y^2}.
\end{equation*}

\section{Реализация}

\subsection{Язык программирования и среда}
\begin{itemize}
    \item \textbf{Используемый язык:} Python 3.14.2
    \item \textbf{Используемая среда разработки:} Visual Studio Code
\end{itemize}

\subsection{Используемые библиотеки}
\begin{itemize}
    \item \textbf{numpy} --- применяется для векторных вычислений, операций с ковариационными матрицами, нахождения собственных значений и векторов, вычисления медиан, а также расчета выборочного коэффициента корреляции Пирсона.
    \item \textbf{scipy.stats} --- используется для генерации многомерных случайных величин (\texttt{multivariate\_normal}), вычисления ранговой корреляции Спирмена и определения квантилей распределения $\chi^2$.
    \item \textbf{math} --- предоставляет тригонометрические функции (\texttt{atan2}, \texttt{degrees}) для аналитического вычисления угла наклона главных осей эллипса рассеивания.
    \item \textbf{matplotlib.pyplot} --- обеспечивает построение диаграмм рассеяния и визуализацию эллипсов равновероятности на координатной плоскости.
    \item \textbf{pathlib} --- используется для определения пути и сохранения итоговых файлов таблиц и графиков.
\end{itemize}

\subsection{Суть реализованных методов}
\begin{itemize}
    \item \texttt{generate\_bivariate\_normal(n, rho, rng)} и \texttt{generate\_mixture(n, rng)} --- функции формирования двумерных выборок. Первая генерирует данные из нормального распределения с использованием заданной ковариационной матрицы, зависящей от параметра $\rho$. Вторая формирует смесь как линейную комбинацию двух независимых выборок с весами $0.9$ и $0.1$.
    \item \texttt{pearson\_corr(x, y)}, \texttt{spearman\_corr(x, y)} и \texttt{quadrant\_corr(x, y)} --- реализуют расчет исследуемых коэффициентов корреляции. Квадрантный коэффициент вычисляется алгоритмически.
    \item \texttt{ellipse\_parameters(x, y, prob)} --- вычисляет полуоси $a$, $b$ и угол наклона $\alpha$ для эллипса рассеивания.
    \item \texttt{run\_experiments()} --- запускает цикл из 1000 итераций генерации выборок заданного объема для вычисления математического ожидания $E(z)$ и дисперсии $D(z)$ точечных оценок коэффициентов корреляции.
    \item \texttt{plot\_normals\_grouped\_by\_sample(...)} и \\ \texttt{plot\_mixture\_grouped\_by\_sample(...)} --- выполняют компоновку диаграмм рассеяния с наложением вычисленных эллипсов в формате группировки по объему выборки.
    \item \texttt{save\_csv\_stats(rows)} и \texttt{save\_csv\_ellipses(rows)} --- обеспечивают объединение и экспорт рассчитанных числовых характеристик и геометрических параметров эллипсов в CSV-файлы.
\end{itemize}

\section{Результаты}

В разделе представлены практические результаты. Оценки коэффициентов корреляции Пирсона ($r$), Спирмена ($r_S$) и квадрантного коэффициента ($r_Q$) получены путем усреднения по 1000 независимым экспериментам для каждого объема выборки $n = 20, 60$ и $100$.

\begin{table}[H]
    \centering
    \caption{Оценки коэффициентов корреляции для двумерного нормального распределения}
    \label{tab:corr_normal}
    \begin{tabular}{|l|c|c|c|c|}
        \hline
        $\mathbf{\rho}$ & \textbf{Коэффициент} & $\mathbf{n=20}$ & $\mathbf{n=60}$ & $\mathbf{n=100}$ \\ \hline
        \multirow{3}{*}{$0.0$} 
        & Пирсона $r$   & $0.0026 \pm 0.2307$ & $-0.0019 \pm 0.1315$ & $0.0009 \pm 0.1012$ \\ \cline{2-5}
        & Спирмена $r_S$ & $0.0044 \pm 0.2322$ & $-0.0024 \pm 0.1324$ & $0.0011 \pm 0.1018$ \\ \cline{2-5}
        & Квадрантный $r_Q$ & $-0.0006 \pm 0.2281$ & $-0.0049 \pm 0.1334$ & $0.0044 \pm 0.1050$ \\ \hline
        \multirow{3}{*}{$0.5$} 
        & Пирсона $r$   & $0.4791 \pm 0.1852$ & $0.4973 \pm 0.0991$ & $0.4966 \pm 0.0759$ \\ \cline{2-5}
        & Спирмена $r_S$ & $0.4465 \pm 0.1943$ & $0.4751 \pm 0.1058$ & $0.4777 \pm 0.0811$ \\ \cline{2-5}
        & Квадрантный $r_Q$ & $0.3038 \pm 0.2203$ & $0.3303 \pm 0.1238$ & $0.3301 \pm 0.0970$ \\ \hline
        \multirow{3}{*}{$0.9$} 
        & Пирсона $r$   & $0.8933 \pm 0.0519$ & $0.8975 \pm 0.0267$ & $0.8990 \pm 0.0195$ \\ \cline{2-5}
        & Спирмена $r_S$ & $0.8639 \pm 0.0708$ & $0.8818 \pm 0.0344$ & $0.8856 \pm 0.0254$ \\ \cline{2-5}
        & Квадрантный $r_Q$ & $0.7014 \pm 0.1668$ & $0.6999 \pm 0.0898$ & $0.7096 \pm 0.0710$ \\ \hline
    \end{tabular}
\end{table}

\begin{table}[H]
    \centering
    \caption{Оценки коэффициентов корреляции для смеси нормальных распределений}
    \label{tab:corr_mixture}
    \begin{tabular}{|l|c|c|c|}
        \hline
        \textbf{Коэффициент} & $\mathbf{n=20}$ & $\mathbf{n=60}$ & $\mathbf{n=100}$ \\ \hline
        Пирсона $r$   & $0.6920 \pm 0.0160$ & $0.6978 \pm 0.0044$ & $0.7002 \pm 0.0027$ \\ \hline
        Спирмена $r_S$ & $0.6575 \pm 0.0197$ & $0.6752 \pm 0.0054$ & $0.6796 \pm 0.0036$ \\ \hline
        Квадрантный $r_Q$ & $0.4850 \pm 0.0395$ & $0.4896 \pm 0.0127$ & $0.4932 \pm 0.0077$ \\ \hline
    \end{tabular}
\end{table}

\begin{table}[H]
    \centering
    \caption{Теоретическая корреляция смеси как линейной комбинации}
    \label{tab:mixture_theory}
    \begin{tabular}{|c|c|c|c|}
        \hline
        $D_X$ & $D_Y$ & $\mathrm{cov}(X,Y)$ & $r_{XY}$ \\ \hline
        $0.9100$ & $0.9100$ & $0.6390$ & $0.7022$ \\ \hline
    \end{tabular}
\end{table}

Геометрические параметры эллипсов рассеивания, усредненные по 1000 независимым экспериментам (большая полуось $a$, малая полуось $b$ и угол наклона главных осей $\alpha$), представлены в таблице \ref{tab:ellipses} для всех рассматриваемых объемов выборки.

\begin{table}[H]
    \centering
    \caption{Усредненные параметры эллипсов рассеивания}
    \label{tab:ellipses}
    \begin{tabular}{|l|c|c|c|c|}
        \hline
        \textbf{Распределение} & \textbf{Параметр} & $\mathbf{n=20}$ & $\mathbf{n=60}$ & $\mathbf{n=100}$ \\ \hline
        \multirow{5}{*}{Нормальное, $\rho=0.0$}
        & $a$ & $2.7706$ & $2.6449$ & $2.5883$ \\ \cline{2-5}
        & $b$ & $2.0490$ & $2.2370$ & $2.2789$ \\ \cline{2-5}
        & $\alpha,\,^\circ$ & $0.6766$ & $1.7883$ & $0.5459$ \\ \cline{2-5}
        & $\bar{x}$ & $-0.0028$ & $-0.0008$ & $0.0004$ \\ \cline{2-5}
        & $\bar{y}$ & $0.0117$ & $-0.0097$ & $0.0038$ \\ \hline

        \multirow{5}{*}{Нормальное, $\rho=0.5$}
        & $a$ & $2.9950$ & $3.0006$ & $2.9952$ \\ \cline{2-5}
        & $b$ & $1.6489$ & $1.7005$ & $1.7149$ \\ \cline{2-5}
        & $\alpha,\,^\circ$ & $43.7322$ & $44.6956$ & $45.0130$ \\ \cline{2-5}
        & $\bar{x}$ & $0.0011$ & $0.0045$ & $0.0035$ \\ \cline{2-5}
        & $\bar{y}$ & $0.0099$ & $0.0034$ & $0.0030$ \\ \hline

        \multirow{5}{*}{Нормальное, $\rho=0.9$}
        & $a$ & $3.3371$ & $3.3595$ & $3.3597$ \\ \cline{2-5}
        & $b$ & $0.7470$ & $0.7665$ & $0.7670$ \\ \cline{2-5}
        & $\alpha,\,^\circ$ & $44.9354$ & $45.0260$ & $45.0350$ \\ \cline{2-5}
        & $\bar{x}$ & $-0.0042$ & $0.0021$ & $-0.0043$ \\ \cline{2-5}
        & $\bar{y}$ & $-0.0010$ & $0.0025$ & $-0.0020$ \\ \hline

        \multirow{5}{*}{Смесь распределений}
        & $a$ & $3.0280$ & $3.0295$ & $3.0448$ \\ \cline{2-5}
        & $b$ & $1.2145$ & $1.2558$ & $1.2646$ \\ \cline{2-5}
        & $\alpha,\,^\circ$ & $45.0555$ & $45.0781$ & $44.9271$ \\ \cline{2-5}
        & $\bar{x}$ & $-0.0051$ & $-0.0019$ & $0.0011$ \\ \cline{2-5}
        & $\bar{y}$ & $-0.0035$ & $0.0027$ & $-0.0018$ \\ \hline
    \end{tabular}
\end{table}

Построены диаграммы рассеяния. Графики нормального закона (рис. 1–3), графики для смеси (рис. 4).

\begin{figure}[H]
    \centering
    \includegraphics[width=1\textwidth]{../figures/normal_n20_by_rho_scatter.png}
    \caption{Диаграммы рассеяния для нормального распределения, $n=20$}
    \label{fig:norm_20}
\end{figure}

\begin{figure}[H]
    \centering
    \includegraphics[width=1\textwidth]{../figures/normal_n60_by_rho_scatter.png}
    \caption{Диаграммы рассеяния для нормального распределения, $n=60$}
    \label{fig:norm_60}
\end{figure}

\begin{figure}[H]
    \centering
    \includegraphics[width=1\textwidth]{../figures/normal_n100_by_rho_scatter.png}
    \caption{Диаграммы рассеяния для нормального распределения, $n=100$}
    \label{fig:norm_100}
\end{figure}

\begin{figure}[H]
    \centering
    \includegraphics[width=1\textwidth]{../figures/mixture_by_n_scatter.png}
    \caption{Диаграммы рассеяния для смеси распределений}
    \label{fig:mixture}
\end{figure}

\FloatBarrier


\section{Обсуждение}

Для двумерного нормального закона выборочный коэффициент Пирсона $r$ выступает наиболее эффективной статистической оценкой линейной зависимости. Эмпирические данные подтверждают минимальный разброс этой величины по сравнению с ранговыми аналогами: при $\rho=0.9$ и $n=100$ стандартное отклонение $\sqrt{D(r)}$ составляет $0.0195$, тогда как для коэффициента Спирмена $\sqrt{D(r_S)} = 0.0254$, а для квадрантного $\sqrt{D(r_Q)} = 0.0710$. Увеличение объема выборки от $n=20$ до $n=100$ приводит к закономерному убыванию дисперсий всех рассматриваемых оценок (табл. \ref{tab:corr_normal}).

Для смеси получено теоретическое значение $r_{XY}\approx 0.7022$ при $D_X=D_Y=0.91$ и $\mathrm{cov}(X,Y)=0.639$ (табл. \ref{tab:mixture_theory}).

Исследование смеси распределений показывает, что коэффициент Пирсона остается устойчиво положительным и быстро приближается к теоретическому значению из табл. \ref{tab:mixture_theory}. При $n=100$ получено $E(r)=0.7002$, а отклонение от теоретического значения $0.7022$ является минимальным среди рассмотренных объемов выборки. Коэффициенты Спирмена и квадрантный также остаются положительными, но дают более консервативные оценки зависимости: $E(r_S)=0.6796$ и $E(r_Q)=0.4932$ (табл. \ref{tab:corr_mixture}).


Для независимых нормальных величин ($\rho=0.0$) полуоси близки ($a=2.5883$, $b=2.2789$ при $n=100$), а угол наклона главной оси остается близким к нулю, что соответствует отсутствию выраженной направленности эллипса. При росте корреляции до $\rho=0.9$ эллипс заметно вытягивается, а угол стабилизируется около $45^\circ$ ($a=3.3597$, $b=0.7670$, $\alpha=45.0350$). Для смеси при $n=100$ получено $a=3.0448$, $b=1.2646$, $\alpha=44.9271$ (табл. \ref{tab:ellipses}).

\section{Выводы}

Для двумерного нормального распределения наиболее эффективной мерой линейной зависимости выступает коэффициент Пирсона $r$. При $\rho=0.9$ и $n=100$ он обеспечивает минимальную дисперсию оценки ($\sqrt{D(r)} = 0.0195$) по сравнению с ранговым $r_S$ и квадрантным $r_Q$ коэффициентами. Увеличение объема выборки закономерно снижает статистический разброс всех вычисляемых характеристик.

Для смеси, реализованной как линейная комбинация двух выборок, выборочный коэффициент Пирсона дает значение, близкое к теоретическому: при $n=100$ получено $E(r)=0.7002$ против теоретического $0.7022$. Ранговый и квадрантный коэффициенты также сохраняют положительный знак, но по величине отстают от Пирсона: $E(r_S)=0.6796$, $E(r_Q)=0.4932$.


Геометрические параметры эллипсов рассеивания жестко зависят от ковариационной структуры. Для независимых случайных величин ($\rho=0.0$) эллипс близок к окружности. Рост коэффициента корреляции до $\rho=0.9$ приводит к выраженному вытягиванию фигуры вдоль главной оси ($a/b \approx 4.38$ для $n=100$), при этом угол наклона стабилизируется около $45^\circ$. Для смеси при $n=100$ получены параметры $a=3.0448$, $b=1.2646$, $\alpha=44.9271$, что соответствует выраженному вытягиванию эллипса и наклону главной оси около $45^\circ$.

\refstepcounter{section}
\addcontentsline{toc}{section}{\thesection\ Список литературы}
\renewcommand{\refname}{\thesection\ Список литературы}

\begin{thebibliography}{9}

    \bibitem{babakhina}
    Бабахина С. А., Баженов А. Н. Практикум по математической статистике: учеб. пособие. — СПб. : Санкт-Петербургский политехнический университет Петра Великого, 2026. 

    \bibitem{maksimov}
    Максимов Ю. Д. Математика. Теория и практика по математической статистике. Конспект-справочник по теории вероятностей : учеб. пособие. — СПб. : Изд-во Политехн. ун-та, 2009. — 395 с. 

\end{thebibliography}

\section{Приложение}

Исходный код программ доступен в репозитории на GitHub: \\
\url{https://github.com/Reichenau/math-stats-labs}
 
\end{document}
